%% =============================================================================
%% sm_d7_design_effects.tex — SM-D, Subsection D.7: Design Effect Diagnostics
%% Body content only (\subsection level); parent structure in sm_d.tex
%% Source: ST4_design_effect_ratios (R-generated), SF6_sandwich_detail.pdf
%% Cross-refs to SM-C DER theory: smc:der, smc:prop-der-classification
%% =============================================================================

\subsection{Design effect diagnostics}
\label{smd:design-effects}

The design effect ratio (DER), defined in \eqref{eq:der} and
classified in~\cref{smc:prop-der-classification}, quantifies the
factor by which the survey design inflates each parameter's variance
beyond the model-based (SRS) variance.  \Cref{smd:tab-der} presents
the full variance decomposition for all 11 fixed-effect parameters,
and~\cref{smd:fig-sandwich} provides a three-panel diagnostic
visualization.

\paragraph{Variance decomposition.}
The three variance quantities in~\cref{smd:tab-der} deserve careful
interpretation.
$[\mathbf{H}_{\mathrm{obs}}^{-1}]_{pp}$ is the model-based variance
obtained from the observed information matrix---the appropriate
benchmark under simple random sampling.
$[\mathbf{V}_{\mathrm{sand}}]_{pp}$ is the sandwich-corrected variance
\citep{Binder1983} that accounts for the complex survey design, and constitutes the
primary inferential summary.
$[\bSigma_{\mathrm{MCMC}}]_{pp}$ is the raw MCMC posterior variance,
which is substantially inflated (by factors of 600--6{,}500) relative
to $\mathbf{H}_{\mathrm{obs}}^{-1}$ because the hierarchical prior
dominates the pseudo-likelihood contribution for fixed effects.

The DER ranges from 1.14 ($\beta_{\mathrm{Hisp}}$) to 4.18
($\alpha_{\mathrm{Urban}}$), with a mean of 2.11.  This is broadly
consistent with the Kish $\DEFF = 3.76$: the DER is systematically
smaller because the state-varying coefficients partially absorb
within-state clustering (\cref{smc:prop-svc-absorption}).  The
corresponding confidence interval inflation factors are
$\sqrt{\DER} \in [1.07, 2.04]$, so the widest Wald intervals are
approximately twice those of a hypothetical SRS design.

\paragraph{Diagnostic figure.}
\Cref{smd:fig-sandwich} decomposes the sandwich correction into three
components.  Panel~(a) shows the hierarchical variance ratio---the ratio of
naive (MCMC-based) to Wald (sandwich-based) confidence interval
widths---on a log scale; all parameters exceed a ratio of~10, and
several exceed~100, confirming that the raw MCMC posterior is
uninformative about the design-consistent uncertainty.  Panel~(b)
compares the DER (based on $\mathbf{H}_{\mathrm{obs}}^{-1}$) with the
``DER\textsubscript{MCMC}'' (based on $\bSigma_{\mathrm{MCMC}}$); the
latter is near zero for all parameters except $\log\kappa$, further
illustrating hierarchical inflation.  Panel~(c) reports the diagonal elements
of the Cholesky scaling matrix $\mathbf{A}$ (\cref{thm:cholesky})
for the fixed-effects block.
For the $\alpha$ and $\beta$ coefficients, the diagonal entries are
all below $0.06$, indicating substantial contraction of the MCMC draws, a direct consequence of the
hierarchically inflated posterior being much wider than the design-consistent
target.

\begin{table}[htbp]
\centering
\footnotesize
\begin{adjustbox}{max width=\textwidth}
\begin{tabular}{@{}l rrr rr rrr@{}}
  \toprule
  Parameter &
  $H_{\mathrm{obs}}^{-1}$ &
  $V_{\mathrm{sand}}$ &
  $\Sigma_{\mathrm{MCMC}}$ &
  DER &
  DER$_{\text{mc}}$ &
  SD($H$) &
  SD(Sand) &
  SD(MCMC) \\
  \midrule
  $\alpha_{\mathrm{Int}}$ & 0.00097 & 0.00198 & 0.946 & 2.04 & 0.002 & 0.031 & 0.044 & 0.972 \\
  $\alpha_{\mathrm{Pov}}$ & 0.00114 & 0.00275 & 0.900 & 2.41 & 0.003 & 0.034 & 0.052 & 0.949 \\
  $\alpha_{\mathrm{Urb}}$ & 0.00028 & 0.00118 & 0.907 & 4.18 & 0.001 & 0.017 & 0.034 & 0.952 \\
  $\alpha_{\mathrm{Blk}}$ & 0.00127 & 0.00195 & 0.947 & 1.54 & 0.002 & 0.036 & 0.044 & 0.973 \\
  $\alpha_{\mathrm{Hisp}}$ & 0.00146 & 0.00265 & 0.920 & 1.82 & 0.003 & 0.038 & 0.051 & 0.959 \\
  $\beta_{\mathrm{Int}}$ & 0.00017 & 0.00023 & 0.860 & 1.36 & 0.000 & 0.013 & 0.015 & 0.927 \\
  $\beta_{\mathrm{Pov}}$ & 0.00022 & 0.00036 & 0.875 & 1.61 & 0.000 & 0.015 & 0.019 & 0.936 \\
  $\beta_{\mathrm{Urb}}$ & 0.00013 & 0.00042 & 0.873 & 3.18 & 0.000 & 0.012 & 0.021 & 0.934 \\
  $\beta_{\mathrm{Blk}}$ & 0.00024 & 0.00033 & 0.845 & 1.38 & 0.000 & 0.016 & 0.018 & 0.919 \\
  $\beta_{\mathrm{Hisp}}$ & 0.00025 & 0.00028 & 0.892 & 1.14 & 0.000 & 0.016 & 0.017 & 0.945 \\
  $\log\kappa$ & 0.00033 & 0.00086 & 0.001 & 2.58 & 1.119 & 0.018 & 0.029 & 0.028 \\
  \bottomrule
\end{tabular}
\end{adjustbox}
\caption{Design effect ratio diagnostics for all 11 fixed-effect
  parameters.  Columns: $H_{\mathrm{obs}}^{-1}$ = model-based variance
  (observed information inverse), $V_{\mathrm{sand}}$ = sandwich
  variance, $\Sigma_{\mathrm{MCMC}}$ = raw MCMC posterior variance,
  DER = $V_{\mathrm{sand}} / H_{\mathrm{obs}}^{-1}$,
  DER$_{\text{mc}}$ = $V_{\mathrm{sand}} / \Sigma_{\mathrm{MCMC}}$,
  and the corresponding standard deviations.  The DER range 1.14--4.18
  (mean 2.11) is consistent with the Kish $\DEFF = 3.76$.}
\label{smd:tab-der}
\end{table}

\begin{figure}[htbp]
  \centering
  \includegraphics[width=\textwidth,height=0.75\textheight,keepaspectratio]{SF6_sandwich_detail.pdf}
  \caption{Sandwich correction diagnostic details.
    (a)~Hierarchical variance ratio (log scale): ratio of naive (MCMC-based)
    to Wald (sandwich-based) confidence interval widths.  Values above
    the dashed line at~10 indicate strong hierarchical inflation.
    (b)~DER comparison: design effect ratio computed using
    $\mathbf{H}_{\mathrm{obs}}^{-1}$ (primary) versus
    $\bSigma_{\mathrm{MCMC}}$ (diagnostic only).
    (c)~Cholesky $\mathbf{A}$ diagonal: scaling factors from the
    affine transformation (\cref{thm:cholesky}).  Values $\ll 1$
    indicate that the transformation substantially contracts the MCMC
    draws toward the design-consistent target.}
  \label{smd:fig-sandwich}
\end{figure}
